\usepackage{changepage}
\usepackage{tikz}
\newcounter{question}
\usepackage{varwidth}
\usepackage{ulem}
\usepackage{multicol}
\usepackage[most]{tcolorbox}
\renewcommand{\ULthickness}{0.5pt}
\setlength{\ULdepth}{1.8pt}
\usepackage{contour}

\contourlength{1pt}
\contourlength{0.8pt}

\newcommand{\myuline}[1]{
  \hspace{-0.5em}
  \uline{\phantom{#1}}%
  \llap{\contour{white}{#1}}%
}
\newenvironment{poemblock}
{%
\begin{varwidth}[t]{0.5\linewidth}
\setlength{\parindent}{1pt}
\setlength{\parskip}{0pt}
}
{%
\end{varwidth}
}

\newcommand{\q}[1]{
\par
\stepcounter{question}
\noindent
\textbf{Câu \thequestion:} #1 
}

\usepackage{calc} % cần cho widthof

\newcommand{\err}[2]{%
\begingroup
\setbox0=\hbox{#1\strut}% ⚠️ đo text gốc, KHÔNG uline
\leavevmode
\myuline{#1}% hiển thị
\rlap{%
  \hspace*{-\wd0}%
  \hspace*{0.5\wd0}%
  \raisebox{-1.2em}{\makebox[0pt]{#2}}%
}%
\endgroup
}
\newcommand{\errorq}[9]{
\par
\stepcounter{question} % <<< THÊM DÒNG NÀY
\noindent
\textbf{Câu \thequestion: }%
#1\,\err{#2}{\textbf{A}}%
\,#3\,\err{#4}{\textbf{B}}%
\,#5\,\err{#6}{\textbf{C}}%
\,#7\,\err{#8}{\textbf{D}}%
\,#9
\par \\
}

\newcommand{\pimg}[2]{{
\par
\noindent
\begin{minipage}{0.68\textwidth}
\setlength{\parindent}{1.5em}
#1
\end{minipage}
\hfill
\begin{minipage}{0.3\textwidth}
\centering
#2
\end{minipage}
\par
}}

\newcommand{\p}[1]{{
\par
\hangindent=0pt
\hangafter=1
\setlength{\parindent}{1.5em}
#1
\par
}}

\newcommand{\source}[2]{%
\par\nobreak
\noindent\makebox[\linewidth][r]{(#1, \textit{#2})}%
\vspace{-1.5em}%
}

\newcommand{\choiceOne}[4]{
\hspace{5em}
\par\noindent
\begin{itemize}[label=, leftmargin=\choicemargin, itemsep=1pt, topsep=0pt]
\item \textbf{A.} #1
\item \textbf{B.} #2
\item \textbf{C.} #3
\item \textbf{D.} #4
\end{itemize}
}

\newcommand{\choiceTwo}[4]{
\hspace{5em}
\par\noindent
\begin{adjustwidth}{\choicemargin}{0pt}
\setlength{\tabcolsep}{0pt}
\renewcommand{\arraystretch}{1}
\begin{tabular}{@{}p{0.48\textwidth}@{\hspace{0.0\textwidth}}p{0.48\textwidth}@{}}
\textbf{A.} #1 & \textbf{B.} #2 \\
\textbf{C.} #3 & \textbf{D.} #4 \\
\end{tabular}
\end{adjustwidth}
}
\newcommand{\choiceFour}[4]{
\hspace{5em}
\par\noindent
\begin{adjustwidth}{\choicemargin}{0pt}
\setlength{\tabcolsep}{0pt}
\renewcommand{\arraystretch}{1}
\begin{tabular}{@{}p{0.24\textwidth}@{\hspace{0.0\textwidth}}
                p{0.24\textwidth}@{\hspace{0.0\textwidth}}
                p{0.24\textwidth}@{\hspace{0.0\textwidth}}
                p{0.24\textwidth}@{}}
\textbf{A.} #1 &
\textbf{B.} #2 &
\textbf{C.} #3 &
\textbf{D.} #4
\end{tabular}
\end{adjustwidth}
}

\newcommand{\reading}[2]{ 
\par
\textbf{Dựa vào thông tin dưới đây để trả lời các câu từ 
\the\numexpr\value{question}+1\relax\ đến 
\the\numexpr\value{question}+#1\relax}
\par
\noindent
#2 
}
\newcommand{\engread}[2]{
\par
\noindent\textbf{Questions 
\the\numexpr\value{question}+1\relax-\the\numexpr\value{question}+#1\relax:} \textit{Read the passage carefully.}%
\par
#2
\par
\noindent
\textit{Choose an option (A, B, C, or D) that best answers each question.}
\par
}
